\clearpage
\setcounter{page}{1}

\section*{Appendix}
\appendix

\noindent In this supplementary material, we provide more details about the RP- and RJ-prompts along with the corresponding descriptions output by the LLMs.


\section{RP- and RJ-prompt}
\label{sec:prompt}

According to the official API interfaces provided by GPT-4 and Llama3, we can interact with large language models (LLMs) in three different roles: ``system'', ``assistant'', and ``user''.
The system message defines the behavior of the assistant within a given context.
The assistant represents the responses and actions of the LLM, while the user refers to the individual engagement with the LLM, providing requests or comments for the assistant to address.
Inspired by the chain-of-thought~\citep{wei2022chain} technique, we have designed a meaningful dialogue as the input prompt for LLMs.
Specifically, the prompts designed for GPT-4 and Llama3-8B are presented below.

\subsection{Prompt for GPT-4}
\label{sec:prompt_gpt}

\noindent \textbf{RP-prompt}.
Relation proposer prompt, which we used to stimulate LLMs to propose all possible relations between two given objects.

\underline{system}: \textit{You are asked to play the role of a} \textbf{relation proposer}\textit{. Given the category names of two objects in an image, you are to infer what kind of relation might exist between them based on your knowledge, and provide the reasons for each possible relation. In the relation between the two objects in the image, we refer to one object as the subject and the other as the object. There may or may not be a relation between the subject and the object. Please note that this relation has an order, that is, the subject comes first and the object comes after. If there is a relation between the two, these relations must belong to one of the pre-defined 56 different types.}

\underline{assistant}: \textit{What are the 56 relations?}

\underline{user}: \textit{They are `over', `in front of', `beside', `on', `in', `attached to', `hanging from', `on back of', `falling off', `going down', `painted on', `walking on', `running on', `crossing', `standing on', `lying on', `sitting on', `flying over', `jumping over', `jumping from', `wearing', `holding', `carrying', `looking at', `guiding', `kissing', `eating', `drinking', `feeding', `biting', `catching', `picking', `playing with', `chasing', `climbing', `cleaning', `playing', `touching', `pushing', `pulling', `opening', `cooking', `talking to', `throwing', `slicing', `driving', `riding', `parked on', `driving on', `about to hit', `kicking', `swinging', `entering', `exiting', `enclosing', `leaning on'.}

\underline{assistant}: \textit{Can you give me an example?}

\underline{user}: \textit{For example, the subject is a person, and the object is a sports ball. The possible relations between them could be: 1. Beside: The person could be standing beside the sports ball. 2. Looking at: The person might be looking at the ball to better control it. 3. Playing: This is because it's very common in real life for a person to be playing with a sports ball. 4. Chasing: The person might be chasing after the ball.}

\underline{assistant}: \textit{Ok, I got it. Please give me the subject and object of the image.}

\underline{user}: \textit{The subject is a SUBJECT\_NAME, and the object is a OBJECT\_NAME.}


\noindent \textbf{RJ-prompt}.
Relation judger prompt, which we used to stimulate LLMs to judge whether there is a specific relation between two objects.

\underline{system}: \textit{You are asked to play the role of a} \textbf{relation judger}\textit{. Given the category names of two objects in an image, and providing you with a relation category name, you need to predict whether this relation is likely to exist in the image based on your knowledge, and give the reason for its existence. For two objects, we call the first object subject and the second object object.}

\underline{assistant}: \textit{Yes, I understand. Can you give me an example?}

\underline{user}: \textit{For example, the input is: the subject is a 'person', the object is a 'sports ball' and the relation is 'playing'. The output should be Yes, the relation is likely to exist in the image. This is because it's very common in real life for a person to be playing with a sports ball.}

\underline{assistant}: \textit{Ok, I got it. Please give me the subject, object and relation names.}

\underline{user}: \textit{The subject is a SUBJECT\_NAME, the object is a OBJECT\_NAME, and the relation is RELATION\_NAME.}


\subsection{Prompt for Llama3-8B}
\label{sec:prompt_llama}
In GPT-4, the ``system'' role can be followed by either the ``assistant'' or the ``user'' role.
However, in Llama3's API, the ``system'' role can only be followed by ``user'' role, not ``assistant'' role.
Therefore, we have made simple modifications to the prompt used for GPT-4. The specific prompt is as follows.

\noindent \textbf{RP-prompt}.
Relation proposer prompt, which we used to stimulate LLMs to propose all possible relations between two given objects.

\underline{system}: \textit{You are asked to play the role of a} \textbf{relation proposer}\textit{. Given the category names of two objects in an image, you are to infer what kind of relation might exist between them based on your knowledge, and provide the reasons for each possible relation. In the relation between the two objects in the image, we refer to one object as the subject and the other as the object. There may or may not be a relation between the subject and the object. Please note that this relation has an order, that is, the subject comes first and the object comes after. If there is a relation between the two, these relations must belong to one of the pre-defined 56 different types. What are the 56 relations?}

\underline{user}: \textit{They are `over', `in front of', `beside', `on', `in', `attached to', `hanging from', `on back of', `falling off', `going down', `painted on', `walking on', `running on', `crossing', `standing on', `lying on', `sitting on', `flying over', `jumping over', `jumping from', `wearing', `holding', `carrying', `looking at', `guiding', `kissing', `eating', `drinking', `feeding', `biting', `catching', `picking', `playing with', `chasing', `climbing', `cleaning', `playing', `touching', `pushing', `pulling', `opening', `cooking', `talking to', `throwing', `slicing', `driving', `riding', `parked on', `driving on', `about to hit', `kicking', `swinging', `entering', `exiting', `enclosing', `leaning on'.}

\underline{assistant}: \textit{Can you give me an example?}

\underline{user}: \textit{For example, the subject is a person, and the object is a sports ball. The possible relations between them could be: 1. Beside: The person could be standing beside the sports ball. 2. Looking at: The person might be looking at the ball to better control it. 3. Playing: This is because it's very common in real life for a person to be playing with a sports ball. 4. Chasing: The person might be chasing after the ball.}

\underline{assistant}: \textit{Ok, I got it. Please give me the subject and object of the image.}

\underline{user}: \textit{The subject is a SUBJECT\_NAME, and the object is a OBJECT\_NAME.}


\noindent \textbf{RJ-prompt}.
Relation judger prompt, which we used to stimulate LLMs to judge whether there is a specific relation between two objects.

\underline{system}: \textit{You are asked to play the role of a} \textbf{relation judger}\textit{. Given the category names of two objects in an image, and providing you with a relation category name, you need to predict whether this relation is likely to exist in the image based on your knowledge, and give the reason for its existence. For two objects, we call the first object subject and the second object object. Please give me an example.}

\underline{user}: \textit{For example, the input is: the subject is a 'person', the object is a 'sports ball' and the relation is 'playing'. The output should be Yes, the relation is likely to exist in the image. This is because it's very common in real life for a person to be playing with a sports ball.}

\underline{assistant}: \textit{Ok, I got it. Please give me the subject, object and relation names.}

\underline{user}: \textit{The subject is a SUBJECT\_NAME, the object is a OBJECT\_NAME, and the relation is RELATION\_NAME.}


\section{RP- and RJ-description}
\label{sec:description}

To gain a more comprehensive understanding of the outputs from LLMs, we selected two sets of subject-object pairs for demonstration, namely ``person-motorcycle'', ``sky-building'', ``person-sea'',  ``rock-horse''.
Here, the former term in each pair represents the category of the subject, and the latter the category of the object.
The RP-descriptions are as below.
For each pair, we randomly select two different relations to showcase their corresponding RJ-descriptions.
Note that in the examples shown below for RP-descriptions, the template sentences that were added to complete the descriptions for all predefined $K$ relations (\eg, it is not likely for SUBJECT\_NAME and OBJECT\_NAME to have the relation RELATION\_NAME) have been omitted. 


\subsection{Description from GPT-4}

\noindent \textbf{RP-description}.
The description generated by LLMs based on the RP-prompt.

\underline{person-motorcycle}: \textit{Based on the given subject (person) and object (motorcycle), here are some possible relations between them: 1.} \textbf{Riding}\textit{: The person could be riding the motorcycle, indicating that they are in control and using it as a means of transportation. 2.} \textbf{Driving}\textit{: Similar to riding, the person could be driving the motorcycle, suggesting that they are operating it. 3.} \textbf{About to hit}\textit{: This relation implies that the person is in a dangerous situation where they are about to collide with the motorcycle, indicating a potential accident or collision. 4.} \textbf{Leaning on}\textit{: The person might be leaning on the motorcycle, using it as a support or resting against it. 5.} \textbf{Parked on}\textit{: The motorcycle could be parked on or near the person, suggesting that they are in close proximity to it but not actively using or interacting with it. 6.} \textbf{Driving on}\textit{: This relation implies that the person is driving another vehicle and the motorcycle is on the road, indicating that they are sharing the same space or path. 7.} \textbf{Looking at}\textit{: The person might be looking at the motorcycle, either out of interest or curiosity, or perhaps considering using it. 8.} \textbf{Touching}\textit{: The person could be touching or interacting with the motorcycle, such as examining it or making adjustments. 9.} \textbf{Holding}\textit{: The person might be holding onto the motorcycle, either for support or to prevent it from falling over. 10.} \textbf{Exiting}\textit{: This relation suggests that the person was previously on or inside the motorcycle and is now in the process of getting off or exiting it. These are just a few possible relations between a person and a motorcycle, and the actual relation would depend on the context and specific details of the image.}

\underline{sky-building}: \textit{Based on the given subject (sky) and object (building), here are some possible relations between them: 1.} \textbf{Over}\textit{: The sky is over the building, as the building is located beneath the sky. 2.} \textbf{In front of}\textit{: The building is in front of the sky, as it is positioned closer to the viewer than the sky. 3.} \textbf{Beside}\textit{: The building is beside the sky, as they are both present in the same scene or landscape. 4.} \textbf{Attached to}\textit{: The building may have structures or elements attached to it that extend into the sky, such as antennas or flagpoles. 5.} \textbf{On}\textit{: The building may have something on its roof or top that is related to the sky, such as a weather vane or satellite dish. 6.} \textbf{In}\textit{: The building may have windows or openings that provide views of the sky from within. 7.} \textbf{Flying over}\textit{: Birds or other flying objects may be seen flying over the building in the sky. 8.} \textbf{About to hit}\textit{: If there is an object in the sky, such as a meteor or a drone, it may appear to be on a collision course with the building. 9.} \textbf{Leaning on}\textit{: If there is a structure or object leaning against the building, it may create a visual relation between the building and the sky. These are just a few possible relations based on general knowledge. The actual relation between the sky and the building in the image may vary depending on the specific context and composition of the scene.}

\underline{person-sea}: \textit{Based on the given subject (person) and object (sea), here are some possible relations between a person and a sea: 1.} \textbf{Standing on}\textit{: The person could be standing on the shore or a pier, with the sea beneath them. 2.} \textbf{Swimming in}\textit{: The person might be swimming in the sea. 3.} \textbf{Looking at}\textit{: The person could be looking at the sea, enjoying the view or observing something in the water. 4.} \textbf{Walking along}\textit{: The person might be walking along the beach or coastline, with the sea beside them. 5.} \textbf{Surfing on}\textit{: The person could be surfing on the waves of the sea. 6.} \textbf{Fishing in}\textit{: The person might be fishing in the sea, either from the shore or from a boat. 7.} \textbf{Diving into}\textit{: The person could be diving into the sea, either for recreational purposes or for activities like snorkeling or scuba diving. 8.} \textbf{Relaxing by}\textit{: The person could be sitting or lying down near the sea, enjoying the peaceful atmosphere. 9.} \textbf{Sailing on}\textit{: The person might be sailing on a boat or yacht on the sea. 10.} \textbf{Collecting seashells}\textit{: The person could be collecting seashells or other items from the shore of the sea. These are just a few possible relations between a person and a sea. The actual relation would depend on the context and specific details of the image.}

\underline{rock-horse}: \textit{Based on the given subject (rock) and object (horse), here are some possible relations between a rock and a horse: 1.} \textbf{Beside}\textit{: The rock could be positioned beside the horse, indicating their proximity to each other. 2.} \textbf{On}\textit{: The horse could be standing on the rock, using it as a platform or a vantage point. 3.} \textbf{In front of}\textit{: The rock could be in front of the horse, blocking its path or serving as an obstacle. 4.} \textbf{Attached to}\textit{: The rock could be attached to the horse, possibly as a weight or burden. These are just a few possible relations between a rock and a horse. The actual relation would depend on the context and specific details of the image.}

\noindent \textbf{RJ-description}.
The description generated by LLMs based on the RJ-prompt.

\underline{person-riding-motorcycle}: \textit{Based on my knowledge, the relation riding is} \textbf{likely} \textit{to exist between person and motorcycle. This is because it is common for a person to ride a motorcycle.}

\underline{person-drinking-motorcycle}: \textit{No, the relation drinking is} \textbf{not likely} \textit{to exist between person and motorcycle. This is because it is not common for a person to be drinking while interacting with a motorcycle.}

\underline{sky-over-building}: \textit{Based on my knowledge, it is} \textbf{likely} \textit{that the relation over exists between a sky and a building in an image. This is because in many real-life scenarios, a sky can be seen positioned above a building.}

\underline{sky-crossing-building}: \textit{No, the relation crossing is} \textbf{not likely} \textit{to exist between sky and building. This is because a sky cannot physically cross a building. The concept of crossing typically involves two objects intersecting or overlapping each other in some way, which is not possible in this case.}

\underline{person-playing-sea}: \textit{No, the relation climbing is} \textbf{not likely} \textit{to exist between person and sea. This is because it is not possible for a person to climb a sea. The relation swimming would be more appropriate in this case.}

\underline{person-sitting on-sea}: \textit{Based on my knowledge, it is} \textbf{unlikely} \textit{for a person to be sitting on a sea. The reason is that a sea is a large body of water and it is not possible for a person to physically sit on it. Therefore, the relation sitting on is unlikely to exist in the image.}

\underline{rock-attached to-horse}: \textit{No, the relation attached to is} \textbf{not likely} \textit{to exist in the image. This is because rocks are typically not attached to horses in real life.}

\underline{rock-falling off-horse}: \textit{No, the relation falling off is} \textbf{not likely} \textit{to exist in the image. This is because rocks cannot fall off horses as they are inanimate objects and do not have the ability to interact with horses in that way.}

\subsection{Description from Llama3-8B}

\noindent \textbf{RP-description}.
The description generated by LLMs based on the RP-prompt.

\underline{person-motorcycle}: \textit{Great! Based on the category names you provided, here are some possible relations between the subject (person) and the object (motorcycle): 1.} \textbf{Riding}\textit{: The person could be riding the motorcycle. 2.} \textbf{Standing next to}\textit{: The person might be standing next to the motorcycle, perhaps admiring it or preparing to ride it. 3.} \textbf{Holding onto}\textit{: The person could be holding onto the motorcycle to steady themselves while riding it. 4.} \textbf{Leaning on}\textit{: The person might be leaning on the motorcycle while taking a break or resting. 5.} \textbf{Fixing}\textit{: The person could be fixing the motorcycle, perhaps making repairs or adjustments. 6.} \textbf{Parked near}\textit{: The person could be parked near the motorcycle, perhaps intending to ride it later. 7.} \textbf{Getting on}\textit{: The person could be getting on the motorcycle, perhaps after having just parked it. 8.} \textbf{Falling off}\textit{: The person might be falling off the motorcycle, perhaps after losing control or experiencing an accident. 9.} \textbf{Facing}\textit{: The person could be facing the motorcycle, perhaps to check on it or to prepare to ride it. 10.} \textbf{Behind}\textit{: The person could be behind the motorcycle, perhaps to push it or to provide support while riding. Please let me know if any of these relations are correct or if you would like me to provide more options.}

\underline{sky-building}: \textit{Great! Based on the category names you provided, here are some possible relations between the subject (sky) and the object (building): 1.} \textbf{Over}\textit{: The building could be located above the sky. 2.} \textbf{Next to}\textit{: The building could be situated next to the sky. 3.} \textbf{Looking at}\textit{: The building might be visible from the sky. 4.} \textbf{Near}\textit{: The building could be near the sky, perhaps located on the same horizon. 5.} \textbf{Part of}\textit{: The building could be part of the sky, perhaps a cloud or a bird flying overhead. 6.} \textbf{Hanging from}\textit{: The building could be suspended from the sky, perhaps by a bridge or a cable. 7.} \textbf{Going through}\textit{: The building could be going through the sky, perhaps as a plane flying through the clouds. 8.} \textbf{Coming from}\textit{: The building could be coming from the sky, perhaps as a hot air balloon rising into the atmosphere. 9.} \textbf{Flying over}\textit{: The building could be flying over the sky, perhaps as a bird or an airplane. 10.} \textbf{Underneath}\textit{: The building could be located underneath the sky, perhaps as a submarine or an underground structure. These are just a few possibilities based on the category names you provided.}

\underline{person-sea}: \textit{Great! Based on the category names you provided, here are some possible relations that could exist between a person and a sea: 1.} \textbf{Beside}\textit{: The person could be standing or sitting beside the sea, enjoying the view or engaging in activities such as swimming, fishing, or boating. 2.} \textbf{Looking at}\textit{: The person might be looking at the sea, perhaps in awe or wonder at its vastness or its role in their life. 3.} \textbf{Swimming}\textit{: The person could be swimming in the sea, either for exercise or recreation. 4.} \textbf{Fishing}\textit{: The person might be fishing in the sea, using various techniques and tools to catch fish or other sea creatures. 5.} \textbf{Boating}\textit{: The person could be on a boat in the sea, either for pleasure or for a specific purpose such as transportation or commercial fishing. 6.} \textbf{Surfing}\textit{: The person might be surfing in the sea, using a surfboard to ride the waves. 7.} \textbf{Sunbathing}\textit{: The person could be lying on the beach or on a boat, soaking up the sun and enjoying the warmth of the sea. 8.} \textbf{Beachcombing}\textit{: The person might be walking along the beach, searching for interesting shells, rocks, or other treasures washed up by the sea. 9.} \textbf{Tidal pooling}\textit{: The person could be exploring the shallow pools of water left behind by the sea at low tide, looking for creatures or interesting rocks. 10.} \textbf{Coastal defense}\textit{: The person might be involved in activities related to protecting the coastline from erosion, pollution, or other threats, such as building seawalls or monitoring water quality. I hope these suggestions help! Let me know if you have any other questions.}

\underline{rock-horse}: \textit{Great! Based on the category names you provided, here are some possible relations between the subject (rock) and the object (horse): 1.} \textbf{Beside}\textit{: The rock could be beside the horse, perhaps as the horse is grazing or resting. 2.} \textbf{Looking at}\textit{: The horse might be looking at the rock, perhaps curious about it or interested in it for some reason. 3.} \textbf{Near}\textit{: The rock could be near the horse, perhaps in a field or meadow where the horse is roaming. 4.} \textbf{Riding}\textit{: The horse could be ridden over the rock, perhaps as a way to traverse a difficult terrain or to reach a destination. 5.} \textbf{Jumping}\textit{: The horse could be jumping over the rock, perhaps in a competition or as a way to clear an obstacle. 6.} \textbf{Kicking}\textit{: The horse could be kicking the rock, perhaps as a way to dislodge it or to mark its territory. 7.} \textbf{Standing on}\textit{: The horse could be standing on the rock, perhaps as a way to get a better view or to rest. 8.} \textbf{Falling off}\textit{: The horse could be falling off the rock, perhaps due to a misstep or a sudden movement. 9.} \textbf{Crossing}\textit{: The horse could be crossing the rock, perhaps as a way to get to the other side of a river or stream. 10.} \textbf{Pushing}\textit{: The horse could be pushing the rock, perhaps as a way to move it out of the way or to create a path. Please let me know if you have any other questions or if you would like me to generate more relations.}

\noindent \textbf{RJ-description}.
The description generated by LLMs based on the RJ-prompt.

\underline{person-riding-motorcycle}: \textit{Based on my knowledge, I would say that the relation riding is} \textbf{likely} \textit{to exist in the image. People often ride motorcycles as a mode of transportation or recreation, so it\'s a common and plausible relation in real life. Therefore, I would give a Yes answer to the question of whether the relation exists in the image.}

\underline{person-drinking-motorcycle}: \textit{Based on my knowledge, the relation drinking is} \textbf{unlikely} \textit{to exist in the image. The reason for this is that it is not common for a person to be drinking a motorcycle. Motorcycles are vehicles that are typically used for transportation, not for drinking. Therefore, the relation drinking does not make sense in this context.}

\underline{sky-over-building}: \textit{Based on the information provided, I would say that the relation over is} \textbf{likely} \textit{to exist in the image. The reason for this is that buildings are often found in the sky, either standing alone or as part of a skyline. The relation over is a common way to describe the position of a building in relation to the sky, as it is typically located above the ground or other objects in the scene. Therefore, given the subject of sky and the object of building, the relation over is a likely candidate for the relation that exists between them in the image.}

\underline{sky-crossing-building}: \textit{Based on my knowledge, the relation crossing is} \textbf{likely} \textit{to exist in the image because the building is located in the sky, and it is common for buildings to be crossed by the sky in many real-life scenarios. Therefore, my answer would be Yes, the relation is likely to exist in the image.}

\underline{person-playing-sea}: \textit{Based on my knowledge, the relation playing is} \textbf{unlikely} \textit{to exist in the image you provided. The subject, person, is not likely to be playing with the object, sea, as they are two very different entities that are not typically associated with each other. The person is a land-dwelling creature, while the sea is a body of water. While it is possible for a person to be near or on the sea, it is not common for them to be playing with it. Therefore, I would say that the relation playing does not likely exist in the image you provided.}

\underline{person-sitting on-sea}: \textit{Based on my knowledge, the relation sitting on is} \textbf{unlikely} \textit{to exist in the image. People do not typically sit on seas, as seas are bodies of water and do not provide a suitable surface for sitting. Therefore, the probability of this relation existing in the image is low.}

\underline{rock-attached to-horse}: \textit{Based on the information provided, I would say that the relation attached to is} \textbf{unlikely} \textit{to exist in the image. A rock and a horse are very different objects and do not have a natural connection or relationship with each other. Therefore, it is unlikely that they would be attached to each other in a real-world scenario.}

\underline{rock-falling off-horse}: \textit{Based on my knowledge, the relation falling off is} \textbf{unlikely} \textit{to exist in an image of a rock and a horse. The reason for this is that rocks and horses are not typically associated with each other in a falling off context. Rocks are inanimate objects that do not have the ability to fall off, while horses are living creatures that are not known to fall off of other objects. Therefore, the probability of seeing a relation of falling off between a rock and a horse in an image is low.}


\section{Description analysis}
\label{sec:analysis_llm}

Through the series of examples above, we analyze the descriptions generated by the LLMs as follows.

\noindent \textbf{Comparation between GPT-4 and Llama3-8B}.
Comparing the descriptions generated by GPT-4 and Llama3-8B, we find that GPT-4's descriptions are slightly better than those of Llama3-8B.
However, thanks to our carefully designed prompts, even using the Llama3-8B model can still produce high-quality descriptions to aid in relation prediction.
As shown in 
Tab.~\ref{tab:ablation_language_feature_extractor} 
% Tab.~4
of our paper, models trained using descriptions from Llama3-8B are only 0.5\% lower in mR@100 compared to those using GPT-4.

\noindent \textbf{RP-descriptions}.
For RP-descriptions, on one hand, some relation predicted by LLMs may occasionally fall out of the predefined set of $K$ relation categories, the mechanism we designed can still handle it though. 
On the other hand, LLMs tend to be more accurate in predicting relations for thing-thing and thing-stuff pairs compared to stuff-thing and stuff-stuff.
For stuff-thing and stuff-stuff pairs, models are more inclined to predict spatial relations.

\noindent \textbf{RJ-descriptions}.
For RJ-descriptions, in most cases, different models provide consistent and reasonable judgments with corresponding explanations.
However, there are occasional inconsistencies, such as with the relation of sky-crossing-building, where GPT-4 deems it impossible, but Llama3-8B considers it possible.
Nevertheless, all LLMs provide reasons to explain the possibility of the relation, and these reasons in the descriptions provide context to the relation judgment, thereby assisting the model in relation prediction.
